\chapter{Asians, Lookbacks, Cliquets and Forward-Starts}\label{ch:dv:asians-lookbacks-cliquets-and-forward-starts}

A corporate treasurer asks for a one-year call on the index, and then for the same call on the average of twelve
monthly prices. The second quote is 40\% cheaper. Nobody on the call can say where the 40\% went, and the treasurer
suspects a trick. There is none. The average of twelve prices moves much less than the last one: its volatility is
about 61\% of the index's. This chapter prices the path-dependent payoffs that desks sell most. Asians average the
path, lookbacks read its extremes, and forward-starts and \omterm{def:dv:asians-lookbacks-cliquets-and-forward-starts:cliquet}{cliquets} are options whose strikes are set in the future.
The first two are priced well by any model that fits today's smile. The last two are not, because their value
depends on the smile the model predicts for future dates. The chapter measures that dependence with two models that
fit the same surface.

\section{Averaging: Asians}

\begin{definition}[Asian option]\label{def:dv:asians-lookbacks-cliquets-and-forward-starts:asian}
An \emph{Asian option}\index{Asian option} pays on the average of the underlying over a set of fixing dates: the
average-rate form pays $(\bar S-K)^+$ (the average-price option of One Quant Book 3, chapter 12), and the average-strike
form pays $(S_T-\bar S)^+$. The average is arithmetic unless stated.
\end{definition}

Averaging lowers volatility. With $n$ fixings equally spaced over $T$, the logarithm of the geometric average is normal
with variance
\[
\sigma_G^2T=\sigma^2T\,\frac{(n+1)(2n+1)}{6n^2},
\]
which tends to $\sigma^2T/3$ as the fixings become continuous. The geometric Asian is therefore a Black--Scholes option
at a reduced volatility and a reduced forward, in closed form. The arithmetic average, the one contracts use, has no
closed form. It is larger than the geometric one, but the two move almost together.

\begin{example}[Where the 40\% went]\label{ex:dv:asians-lookbacks-cliquets-and-forward-starts:asian}
Spot and strike 100, one year, $r=3\%$, $q=1\%$, volatility 20\%. The vanilla call is worth 8.83, and the call on the
average of twelve monthly fixings is worth 5.32, 40\% less. The variance factor is $13\times25/(6\times144)=0.376$, so the
geometric average moves with a volatility of $0.61\times20\%=12.3\%$. As the fixings multiply, the Asian's price relative
to the vanilla falls from 0.78 with two fixings to 0.68 with four, 0.60 with twelve and 0.57 with daily fixings. It tracks
the volatility factor $\sqrt{(n+1)(2n+1)/(6n^2)}$: 0.79, 0.68, 0.61 and 0.58
(\cref{fig:dv:asians-lookbacks-cliquets-and-forward-starts:asian}, left).
\end{example}

The near-perfect co-movement of the two averages is what makes Monte Carlo pricing of Asians cheap. The geometric Asian
is a control variate (One Quant Book 4, chapter 26). Kemna and Vorst proposed it. On each path compute both payoffs, and
correct the arithmetic estimate by the geometric payoff's known error:
$\hat A_{\mathrm{cv}}=\bar A-\beta(\bar G-G_{\mathrm{exact}})$, with $\beta$ estimated by regression on the same
paths. On the example the two discounted payoffs have a correlation of 0.9996 and $\beta=1.03$. The standard error falls
from 0.025 to 0.0007, a variance reduction of 1\,357 times: one path with the control variate is worth 1\,357 without
it (\cref{fig:dv:asians-lookbacks-cliquets-and-forward-starts:asian}, right).

\begin{omfigure}
\begin{tikzpicture}
\begin{groupplot}[group style={group size=2 by 1,horizontal sep=1.6cm},
  width=6.6cm,height=5cm,
  tick label style={font=\small},label style={font=\small},grid=major,grid style={gray!25},table/col sep=comma]
\nextgroupplot[xmode=log,xlabel={fixings in the year},ylabel={ratio},xmin=0.8,xmax=320,ymin=0.5,ymax=1.05,
  xtick={1,2,4,12,52,252},xticklabels={1,2,4,12,52,252},
  title={\small Asian / vanilla},
  legend columns=1,legend cell align=left,legend style={font=\footnotesize,at={(0.5,-0.32)},anchor=north,draw=none,fill=none}]
\addplot[omThm,very thick,mark=*,mark size=1.6pt] table[x=n,y=ratio]{figdata/derivatives/16-asians-lookbacks-cliquets-and-forward-starts/asian_fixings.csv};
\addplot[omDef,thick,dashed] table[x=n,y=volfactor]{figdata/derivatives/16-asians-lookbacks-cliquets-and-forward-starts/asian_fixings.csv};
\legend{price ratio,{$\sqrt{(n+1)(2n+1)/6n^2}$}}
\nextgroupplot[xlabel={geometric-average payoff},ylabel={arithmetic-average payoff},xmin=0,xmax=40,ymin=0,ymax=40,
  title={\small 400 paths, twelve fixings}]
\addplot[omThm,only marks,mark=*,mark size=0.8pt] table[x=geometric,y=arithmetic]{figdata/derivatives/16-asians-lookbacks-cliquets-and-forward-starts/asian_scatter.csv};
\draw[gray,dashed] (axis cs:0,0) -- (axis cs:40,40);
\end{groupplot}
\end{tikzpicture}

\omcaption{Left: the price of an at-the-money one-year Asian call relative to the vanilla, by the number of fixings,
against the volatility factor of the geometric average. Right: the arithmetic and geometric payoffs path by path: nearly
the same random variable, which is why the geometric Asian, known in closed form, is an efficient control variate. Data:
the tutorial.}\label{fig:dv:asians-lookbacks-cliquets-and-forward-starts:asian}
\end{omfigure}

Asians have two practical properties that the price alone does not show. Their risk falls as fixings pass: once half the
fixings are set, half the payoff is known, and the delta and \omterm{def:dv:greeks-and-the-hedging-pnl:greeks}{vega} shrink accordingly. Hedging near a fixing, however,
concentrates flow on the fixing date and time, one reason averages are taken over many fixings. And the smile's
dynamics matter little. An Asian's value depends mostly on the at-the-money region of the smile at each fixing date,
which any model fitted to the surface reproduces. On chapter 9's surface (zero rates, twelve monthly fixings), the
at-the-money Asian is worth 4.34 under \omterm{def:dv:local-volatility:model}{local volatility} and 4.36 under Heston, within one standard error of each other.
Black--Scholes at the one-year at-the-money volatility gives 4.69: what matters is the term structure, since the early
fixings see the lower short-dated volatilities.

\section{Extremes: lookbacks}

\begin{definition}[Lookback option]\label{def:dv:asians-lookbacks-cliquets-and-forward-starts:lookback}
A \emph{lookback option}\index{lookback option} pays on the maximum or minimum of the underlying over its life. The
floating-strike lookback call pays $S_T-\min_tS_t$ (buy at the low); the fixed-strike lookback call pays
$(\max_tS_t-K)^+$.
\end{definition}

A lookback is a barrier option integrated over all barrier levels: $S_T-\min S=\int\mathbf 1\{\min S<h\}\,dh$ plus
the terminal price. Goldman, Sosin and Gatto priced it in closed form with the reflection principle, for continuous
monitoring. It is expensive. With the example's parameters the floating-strike lookback call is worth 15.69, 1.8 times the
at-the-money call, because it buys at the best price of the year.

Contracts observe the extreme at discrete dates, often daily closes, and the monitored minimum is higher than the true
one. The \omterm{def:dv:barriers-and-digitals:shift}{barrier shift} of chapter 15 carries over: the monitored minimum behaves like the continuous one multiplied by
$e^{\beta\sigma\sqrt{\Delta t}}$, which gives $C_{\mathrm d}\approx e^{x}C_{\mathrm c}-(e^x-1)Se^{-qT}$ with
$x=\beta\sigma\sqrt{\Delta t}$. With daily monitoring a simulation gives 15.09 (standard error 0.04) and the shifted formula
15.08. Weekly they give 14.42 and 14.34, monthly 13.13 and 12.84, and with four dates 11.55 and 10.70
(\cref{fig:dv:asians-lookbacks-cliquets-and-forward-starts:lookback}). As with barriers, the correction serves daily
monitoring, and a sparse schedule needs the simulation.

\begin{omfigure}
\begin{tikzpicture}
\begin{axis}[width=10.5cm,height=5cm,xmode=log,
  xlabel={monitoring dates a year},ylabel={floating-strike lookback call},
  tick label style={font=\small},label style={font=\small},
  xmin=3,xmax=320,ymin=8,ymax=16.5,grid=major,grid style={gray!25},
  xtick={4,12,52,252},xticklabels={4,12,52,252},
  legend columns=4,legend style={font=\footnotesize,at={(0.5,-0.3)},anchor=north,draw=none,fill=none,/tikz/every even column/.append style={column sep=6pt}},
  table/col sep=comma]
\addplot[black,only marks,mark=*,mark size=1.8pt] table[x=n,y=mc]{figdata/derivatives/16-asians-lookbacks-cliquets-and-forward-starts/lookback.csv};
\addplot[omThm,very thick] table[x=n,y=shifted]{figdata/derivatives/16-asians-lookbacks-cliquets-and-forward-starts/lookback.csv};
\addplot[gray,thick,dashed] table[x=n,y=cont]{figdata/derivatives/16-asians-lookbacks-cliquets-and-forward-starts/lookback.csv};
\addplot[omDef,thick,dotted] table[x=n,y=vanilla]{figdata/derivatives/16-asians-lookbacks-cliquets-and-forward-starts/lookback.csv};
\legend{Monte Carlo,shifted extreme,continuous,vanilla call}
\end{axis}
\end{tikzpicture}

\omcaption{The floating-strike lookback call (one year, 20\%, $r=3\%$, $q=1\%$) by monitoring frequency: simulation, the
continuous formula of Goldman, Sosin and Gatto, and the same with the extreme shifted as a barrier is. Data: the
tutorial.}\label{fig:dv:asians-lookbacks-cliquets-and-forward-starts:lookback}
\end{omfigure}

A lookback's hedge is a ladder of barrier options. Each time the price makes a new low, the strike resets and the desk's
position in the option below changes. The payoff is continuous, so there is no gap at a single level, but the delta
jumps each time a new extreme is set, and the \omterm{def:dv:greeks-and-the-hedging-pnl:greeks}{vega} is large. It is a volatility product more than a directional one.

\section{Forward-starts and the forward smile}

\begin{definition}[Forward-start option]\label{def:dv:asians-lookbacks-cliquets-and-forward-starts:fwdstart}
A \emph{forward-start option}\index{forward-start option} is an option whose strike is fixed at a future date $t_1$ as a
proportion $k$ of the spot then, $K=kS_{t_1}$, and which expires at $t_2>t_1$. It pays
$(S_{t_2}-kS_{t_1})^+$ for a call.
\end{definition}

Under Black--Scholes the price follows from homogeneity. At $t_1$ the option is a call with strike $kS_{t_1}$, worth
$S_{t_1}C_{\mathrm{BS}}(1,k,t_2-t_1)$, and the expected discounted $S_{t_1}$ is $S_0e^{-qt_1}$, so the price is
$S_0e^{-qt_1}C_{\mathrm{BS}}(1,k,t_2-t_1)$. It depends only on the volatility between $t_1$ and $t_2$ and on the
smile that will prevail at $t_1$. In any other model the implied volatilities of \omterm{def:dv:asians-lookbacks-cliquets-and-forward-starts:fwdstart}{forward-start options} across $k$ form
the \omterm{def:dv:local-volatility:forward}{forward smile} of chapter 9. That smile is not observed today: it is the model's forecast.

\begin{example}[Two forward smiles from one surface]\label{ex:dv:asians-lookbacks-cliquets-and-forward-starts:fwdsmile}
Chapter 9's \omterm{def:dv:local-volatility:model}{local volatility} and chapter 10's \omterm{def:dv:stochastic-volatility:heston}{Heston model} are both fitted to chapter 9's surface: they reprice its
one-year at-the-money call at 19.19\% and 18.99\% against the market's 19.19\%. Their one-month smiles starting in six
months are very different (\cref{fig:dv:asians-lookbacks-cliquets-and-forward-starts:fwdsmile}). \omterm{def:dv:local-volatility:model}{Local volatility} gives
23.9\% at 94\%, 21.5\% at the money and 21.3\% at 106\%, almost flat. Heston gives 24.0\%, 18.5\% and 20.0\%, with a skew
and a smile. Today's one-month smile runs from 19.6\% to 14.9\% to 12.0\%. \omterm{def:dv:local-volatility:model}{Local volatility} flattens the future smile, as
chapter 9 found. Heston keeps its shape, and Bergomi observed that its \omterm{def:dv:local-volatility:forward}{forward smiles} are more convex than today's.
\end{example}

\begin{omfigure}
\begin{tikzpicture}
\begin{axis}[width=10.5cm,height=5cm,
  xlabel={strike (\% of the spot at the start)},ylabel={1-month implied volatility (\%)},
  tick label style={font=\small},label style={font=\small},
  xmin=93,xmax=107,ymin=10,ymax=26,grid=major,grid style={gray!25},
  legend columns=3,legend cell align=left,legend style={font=\footnotesize,at={(0.5,-0.3)},anchor=north,draw=none,fill=none,/tikz/every even column/.append style={column sep=8pt}},
  table/col sep=comma]
\addplot[gray,thick,dashed,mark=o,mark size=1.6pt] table[x=k,y=today]{figdata/derivatives/16-asians-lookbacks-cliquets-and-forward-starts/forward_smiles.csv};
\addplot[omDef,very thick,mark=square*,mark size=1.6pt] table[x=k,y=lv]{figdata/derivatives/16-asians-lookbacks-cliquets-and-forward-starts/forward_smiles.csv};
\addplot[omThm,very thick,mark=*,mark size=1.6pt] table[x=k,y=heston]{figdata/derivatives/16-asians-lookbacks-cliquets-and-forward-starts/forward_smiles.csv};
\legend{today's 1-month smile,local volatility: in 6 months,Heston: in 6 months}
\end{axis}
\end{tikzpicture}

\omcaption{The one-month smile six months from now, as forecast by two models that fit the same surface today, against
today's one-month smile. \omterm{def:dv:local-volatility:model}{Local volatility}'s \omterm{def:dv:local-volatility:forward}{forward smile} is nearly flat; Heston's keeps a skew and a smile. Data: the
tutorial.}\label{fig:dv:asians-lookbacks-cliquets-and-forward-starts:fwdsmile}
\end{omfigure}

Two effects show in the example. The level is the forward volatility: the surface's term structure rises, so the
volatility from six to seven months is above today's one-month volatility in both models. The shape is the model's
dynamics. In \omterm{def:dv:local-volatility:model}{local volatility} the skew is a function of the spot level and time, fixed at calibration, and it decays
with the forward date. In Heston the skew comes from the correlation of the spot with a variance that is still random,
and it is regenerated at every date. The market's own \omterm{def:dv:local-volatility:forward}{forward smile} can be read only from instruments that pay on it:
forward-starts, \omterm{def:dv:asians-lookbacks-cliquets-and-forward-starts:cliquet}{cliquets}, and the forward skew implied by options on the volatility index.

\section{Cliquets}

\begin{definition}[Cliquet]\label{def:dv:asians-lookbacks-cliquets-and-forward-starts:cliquet}
A \emph{cliquet}\index{cliquet} (ratchet) is a series of \omterm{def:dv:asians-lookbacks-cliquets-and-forward-starts:fwdstart}{forward-start options} on consecutive periods. In its common
structured form it pays the sum of the period returns $r_i=S_{t_i}/S_{t_{i-1}}-1$, each clipped between a local floor
and a local cap, with the sum clipped between a global floor and a global cap.
\end{definition}

\begin{definition}[Reverse cliquet]\label{def:dv:asians-lookbacks-cliquets-and-forward-starts:reverse}
A \emph{reverse cliquet}\index{reverse cliquet} pays a large coupon eroded by every negative period return:
$\max\bigl(0,\,C+\sum_i\min(r_i,0)\bigr)$. The investor is short a strip of forward-start puts, capped in total.
\end{definition}

A locally capped and floored period return $\min(\max(r_i,f),c)$ is a forward-start call spread less a forward-start put
spread. When the cap and floor are close to zero it behaves like a forward-start digital, and chapter 15 showed that a
digital's price is set by the skew. A \omterm{def:dv:asians-lookbacks-cliquets-and-forward-starts:cliquet}{cliquet} is therefore a basket of bets on the forward skew, which is exactly the
quantity on which the two models of the last section disagree. The \omterm{def:dv:asians-lookbacks-cliquets-and-forward-starts:reverse}{reverse cliquet} is short forward puts, a bet on the
level of forward volatility and on how fat its left tail is.

\begin{example}[The cliquet that fits every vanilla]\label{ex:dv:asians-lookbacks-cliquets-and-forward-starts:cliquet}
A one-year \omterm{def:dv:asians-lookbacks-cliquets-and-forward-starts:cliquet}{cliquet} with monthly periods pays the sum of the twelve monthly returns, each clipped to $[-1\%,+1\%]$, floored
at zero overall. On 100\,000 paths with zero rates it is worth 1.53\% of notional under \omterm{def:dv:local-volatility:model}{local volatility}, 2.22\% under
Heston and 1.19\% under Black--Scholes at the one-year at-the-money volatility. All three fit the one-year at-the-money
call, and the two smile models fit the whole surface. The Heston price is 45\% above the local-volatility price. The gap
widens as the local cap narrows and the payoff looks more like a digital: at $\pm0.5\%$ the prices are 0.78\% and 1.16\%. A
\omterm{def:dv:asians-lookbacks-cliquets-and-forward-starts:reverse}{reverse cliquet} paying $\max(0,25\%+\sum\min(r_i,0))$ is worth 5.43\% under \omterm{def:dv:local-volatility:model}{local volatility}, 7.56\% under Heston and 3.57\%
under Black--Scholes (\cref{fig:dv:asians-lookbacks-cliquets-and-forward-starts:cliquet}).
\end{example}

\begin{omfigure}
\begin{tikzpicture}
\begin{axis}[width=10.5cm,height=5cm,
  xlabel={local cap and floor ($\pm$, \%)},ylabel={cliquet price (\% of notional)},
  tick label style={font=\small},label style={font=\small},
  xmin=0,xmax=5.2,ymin=0,ymax=6.5,grid=major,grid style={gray!25},
  legend columns=3,legend style={font=\footnotesize,at={(0.5,-0.3)},anchor=north,draw=none,fill=none,/tikz/every even column/.append style={column sep=8pt}},
  table/col sep=comma]
\addplot[omThm,very thick,mark=*,mark size=1.6pt] table[x=cap,y=heston]{figdata/derivatives/16-asians-lookbacks-cliquets-and-forward-starts/cliquet_caps.csv};
\addplot[omDef,very thick,mark=square*,mark size=1.6pt] table[x=cap,y=lv]{figdata/derivatives/16-asians-lookbacks-cliquets-and-forward-starts/cliquet_caps.csv};
\addplot[gray,thick,dashed,mark=o,mark size=1.6pt] table[x=cap,y=bs]{figdata/derivatives/16-asians-lookbacks-cliquets-and-forward-starts/cliquet_caps.csv};
\legend{Heston,local volatility,Black--Scholes}
\end{axis}
\end{tikzpicture}

\omcaption{A one-year monthly \omterm{def:dv:asians-lookbacks-cliquets-and-forward-starts:cliquet}{cliquet} (local cap and floor $\pm c$, global floor zero) under three models that fit the
one-year at-the-money call, and two that fit the whole surface. The narrower the local cap, the more the payoff is a
string of forward digitals and the more the forward skew, not today's smile, sets the price. Data: the
tutorial.}\label{fig:dv:asians-lookbacks-cliquets-and-forward-starts:cliquet}
\end{omfigure}

In 2004 Bergomi counted \omterm{def:dv:asians-lookbacks-cliquets-and-forward-starts:reverse}{reverse cliquets} and Napoleons among the most recent exotic structures, and cited a review that
Risk had run in February of that year under the title ``\omterm{def:dv:asians-lookbacks-cliquets-and-forward-starts:reverse}{Reverse cliquets}: end of the road?''. His argument was that such
products are priced by the dynamics a model assumes for future smiles, not by today's surface.

\section{Which model for which payoff}

The question a structurer asks of a model is not whether it fits the surface. \omterm{def:dv:local-volatility:model}{Local volatility}, Heston with jumps, SABR
by expiry and the \omterm{def:dv:rough-volatility-and-forward-variance-models:bergomi}{Bergomi models} all fit it, or nearly. The question is which quantities beyond the surface the payoff
depends on, and whether the model's assumptions about them are credible.

\begin{center}
\begin{tabular}{p{3.6cm}p{5.2cm}p{4.2cm}}
\toprule
payoff & depends on & adequate models\\
\midrule
European, Asian & terminal distributions at the fixing dates & any model that fits the surface\\
barrier, lookback & the path's extremes: the smile and its dynamics near the barrier & local or local-stochastic volatility\\
forward-start, \omterm{def:dv:asians-lookbacks-cliquets-and-forward-starts:cliquet}{cliquet} & the \omterm{def:dv:local-volatility:forward}{forward smile}, the forward skew & stochastic volatility, Bergomi\\
\omterm{def:dv:asians-lookbacks-cliquets-and-forward-starts:reverse}{reverse cliquet}, Napoleon & forward volatility and its volatility & Bergomi, rough models\\
options on variance or the index & the volatility of variance & \omterm{def:dv:rough-volatility-and-forward-variance-models:fvm}{forward-variance models} (chapter 12)\\
\bottomrule
\end{tabular}
\end{center}

Bergomi's 2004 example is the clearest statement of the principle. A six-year Napoleon paid fixed coupons for two years,
then 8\% a year augmented by the worst of the twelve monthly performances of the Eurostoxx 50 in each year, floored at
zero. Such a product is in effect a put on long-dated forward volatility, whose value a \omterm{def:dv:black-scholes-three-ways:model}{Black--Scholes model} gives
without any charge for the volatility of that volatility. His ``Smile dynamics II'' model was built to price such
products, the Napoleon and the \omterm{def:dv:asians-lookbacks-cliquets-and-forward-starts:reverse}{reverse cliquet} among its examples, with separate control of the short forward skew, the
\omterm{def:dv:stochastic-volatility:sv}{spot--volatility correlation} and the term structure of the \omterm{def:dv:stochastic-volatility:sv}{volatility of volatility}. When two fitted models disagree by
45\%, as in \cref{ex:dv:asians-lookbacks-cliquets-and-forward-starts:cliquet}, the desk prices with the model whose
dynamics it can defend and holds a reserve for the difference (chapter 27).

\section{Tutorial: payoffs on a common path interface}\label{tut:dv:asians-lookbacks-cliquets-and-forward-starts:tut}

\begin{tutorial}
\textbf{Goal.} Price an arithmetic Asian with a geometric control variate, a discretely monitored lookback, and a
forward-start and a \omterm{def:dv:asians-lookbacks-cliquets-and-forward-starts:cliquet}{cliquet} under three models that fit the same at-the-money volatility. \textbf{End state:} the four
figures and the numbers of the weekend problem.
\begin{tutsteps}
\item \textbf{The control variate.} Both payoffs on the same paths, $\beta$ by regression, the geometric error removed:
\omcode{code/firm/pathdep/firm_pathdep.py}{54}{70}{Arithmetic Asian with the geometric control variate.}\label{lst:dv:asians-lookbacks-cliquets-and-forward-starts:cv}
\item \textbf{Payoffs read paths only.} Period returns, the \omterm{def:dv:asians-lookbacks-cliquets-and-forward-starts:cliquet}{cliquet} and the \omterm{def:dv:asians-lookbacks-cliquets-and-forward-starts:reverse}{reverse cliquet}:
\omcode{code/firm/pathdep/firm_pathdep.py}{107}{121}{Cliquet payoffs on the common path array.}\label{lst:dv:asians-lookbacks-cliquets-and-forward-starts:cliquet}
\item \textbf{Paths from each model}: \texttt{firm\_localvol.simulate} with recorded dates, \texttt{firm\_heston.simulate\_paths},
\texttt{gbm\_paths}; then run \texttt{dv\_pathdep.forward\_smiles()}, \texttt{cliquets()} and \texttt{fig\_pathdep.py}.
\end{tutsteps}
\textbf{What to change next.} Price the \omterm{def:dv:asians-lookbacks-cliquets-and-forward-starts:cliquet}{cliquet} with Heston's \omterm{def:dv:stochastic-volatility:sv}{volatility of volatility} halved and refit $\rho$ to keep the
one-year skew; add a global cap of 5\% and see which model's price moves most; price a Napoleon-like coupon, 8\% plus the
worst monthly return, floored at zero.
\end{tutorial}

\section{Build: path-dependent payoffs}\label{bld:dv:asians-lookbacks-cliquets-and-forward-starts:pathdep}

\begin{build}
\bfield{Purpose} The miniature firm's library of path-dependent payoffs, written against one path array so that any
simulator (Black--Scholes, \omterm{def:dv:local-volatility:model}{local volatility}, Heston, and chapter 23's engine) prices any of them.
\bfield{Interface} \texttt{gbm\_paths(s, times, r, q, vol, n\_paths, seed)}; \texttt{asian\_payoff(paths, strike, right,
geometric)}, \texttt{geometric\_asian}, \texttt{asian\_mc}; \texttt{lookback\_payoff}, \texttt{lookback\_floating\_call},
\texttt{lookback\_shifted}; \texttt{forward\_start\_call}; \texttt{period\_returns}, \texttt{cliquet\_payoff(paths,
local\_floor, local\_cap, global\_floor, global\_cap)}, \texttt{reverse\_cliquet\_payoff}; \texttt{forward\_smile\_from\_paths}.
\bfield{Rules} Column 0 of a path array is today's spot; payoffs never simulate; every Monte Carlo price is reported with its
standard error; control variates are used whenever a closed-form cousin exists.
\bfield{Acceptance tests} \texttt{code/firm/pathdep/tests/}: the one-fixing geometric Asian is the vanilla and the formula
matches a simulation; the control variate cuts the error tenfold at least; the lookback exceeds the vanilla and its shifted
formula matches a daily simulation; the forward-start formula and a flat model's \omterm{def:dv:local-volatility:forward}{forward smile}; \omterm{def:dv:asians-lookbacks-cliquets-and-forward-starts:cliquet}{cliquet} payoffs on hand
paths.
\bfield{Stretch} Moment-matching approximations for arithmetic Asians; fixed-strike lookbacks; the Napoleon.
\end{build}

\begin{omsources}
\item A.~G.~Z.~Kemna and A.~C.~F.~Vorst, ``A pricing method for options based on average asset values'', \emph{Journal of
Banking and Finance} 14(1) (1990) 113--129.
\item M.~B.~Goldman, H.~B.~Sosin and M.~A.~Gatto, ``Path dependent options: buy at the low, sell at the high'',
\emph{Journal of Finance} 34(5) (1979) 1111--1127.
\item L.~Bergomi, ``Smile dynamics'', \emph{Risk} (September 2004); ``Smile dynamics II'', SSRN 1493302.
\item C.~Jeffery, ``Reverse cliquets: end of the road?'', \emph{Risk} (February 2004) 20--22.
\end{omsources}

\section{Exercises}

\begin{exercise}[$\star$]\label{exo:dv:asians-lookbacks-cliquets-and-forward-starts:1}
Compute the volatility factor $\sqrt{(n+1)(2n+1)/(6n^2)}$ for $n=4$ and its limit as $n\to\infty$.
\end{exercise}

\begin{exercise}[$\star$]\label{exo:dv:asians-lookbacks-cliquets-and-forward-starts:2}
Why is the arithmetic Asian worth more than the geometric one?
\end{exercise}

\begin{exercise}[$\star$]\label{exo:dv:asians-lookbacks-cliquets-and-forward-starts:3}
Under Black--Scholes with $q=2\%$, what is a six-month-into-six-month at-the-money forward-start call worth relative to a
six-month at-the-money call?
\end{exercise}

\begin{exercise}[$\star\star$]\label{exo:dv:asians-lookbacks-cliquets-and-forward-starts:4}
Show that $\min(\max(r,f),c)=f+(r-f)^+-(r-c)^+$ and explain why a \omterm{def:dv:asians-lookbacks-cliquets-and-forward-starts:cliquet}{cliquet} with a narrow local cap and floor behaves like a
string of forward digitals.
\end{exercise}

\begin{exercise}[$\star\star$]\label{exo:dv:asians-lookbacks-cliquets-and-forward-starts:5}
Explain why the local-volatility \omterm{def:dv:local-volatility:forward}{forward smile} is flatter than Heston's in
\cref{ex:dv:asians-lookbacks-cliquets-and-forward-starts:fwdsmile}.
\end{exercise}

\begin{exercise}[$\star\star$]\label{exo:dv:asians-lookbacks-cliquets-and-forward-starts:6}
A variance-reduction factor of 1\,357 means what, in paths and in computing time, for a price wanted to 0.001?
\end{exercise}

\begin{exercise}[$\star\star\star$]\label{exo:dv:asians-lookbacks-cliquets-and-forward-starts:7}
\emph{Coding.} Price the \omterm{def:dv:asians-lookbacks-cliquets-and-forward-starts:reverse}{reverse cliquet} of \cref{ex:dv:asians-lookbacks-cliquets-and-forward-starts:cliquet} with coupons of
15\% and 35\% under the three models. Which model's price is most sensitive to the coupon?
\end{exercise}

\begin{exercise}[$\star\star\star$]\label{exo:dv:asians-lookbacks-cliquets-and-forward-starts:8}
\emph{Find the flaw.} ``Our local-volatility model reprices every vanilla on the surface to a tenth of a point, so its \omterm{def:dv:asians-lookbacks-cliquets-and-forward-starts:cliquet}{cliquet}
prices are market prices.''
\end{exercise}

\section{Problem: The Cliquet That Fits Every Vanilla}

\begin{problem}[{Weekend problem --- one surface, two prices}]\label{pb:dv:asians-lookbacks-cliquets-and-forward-starts:1}
A structurer must quote a one-year monthly \omterm{def:dv:asians-lookbacks-cliquets-and-forward-starts:cliquet}{cliquet} (local cap and floor $\pm1\%$, global floor zero) on the index of
chapter 9's surface. The desk has two calibrated models: chapter 9's \omterm{def:dv:local-volatility:model}{local volatility} and chapter 10's Heston.

\medskip\noindent\textbf{Part I --- Warm-up: averages and extremes.}
\begin{enumerate}
\item Give the one-year twelve-fixing Asian call and the vanilla (spot 100, 20\%, $r=3\%$, $q=1\%$).
\item Give the variance factor and the volatility of the geometric average.
\item Give the standard errors with and without the control variate.
\item Give the continuous floating-strike lookback call and its daily-monitored value.
\item Why does the shift fail with four monitoring dates?
\end{enumerate}

\medskip\noindent\textbf{Part II --- \omterm{def:dv:local-volatility:forward}{Forward smiles}.}
\begin{enumerate}[resume]
\item Check both models' one-year at-the-money volatility against the market's.
\item Give each model's one-month at-the-money volatility starting in six months, and today's one-month.
\item Give each model's forward skew between 94\% and 106\%, and today's.
\item Why is the forward at-the-money volatility above today's one-month volatility in both models?
\item Which model's \omterm{def:dv:local-volatility:forward}{forward smile} would you expect to match a market for \omterm{def:dv:asians-lookbacks-cliquets-and-forward-starts:fwdstart}{forward-start options}, and why?
\end{enumerate}

\medskip\noindent\textbf{Part III --- The \omterm{def:dv:asians-lookbacks-cliquets-and-forward-starts:cliquet}{cliquet}.}
\begin{enumerate}[resume]
\item Give the \omterm{def:dv:asians-lookbacks-cliquets-and-forward-starts:cliquet}{cliquet}'s price under \omterm{def:dv:local-volatility:model}{local volatility}, Heston and Black--Scholes.
\item Give the prices with a local cap and floor of $\pm0.5\%$ and $\pm5\%$.
\item Why does the gap shrink as the cap widens?
\item Give the \omterm{def:dv:asians-lookbacks-cliquets-and-forward-starts:reverse}{reverse cliquet}'s prices with a 25\% coupon.
\item Which model would a seller of \omterm{def:dv:asians-lookbacks-cliquets-and-forward-starts:reverse}{reverse cliquets} prefer, and why is that a warning?
\end{enumerate}

\medskip\noindent\textbf{Part IV --- Judgement.}
\begin{enumerate}[resume]
\item Which market instruments would tell you which model is right?
\item How would you hedge the \omterm{def:dv:asians-lookbacks-cliquets-and-forward-starts:cliquet}{cliquet}'s forward-skew exposure?
\item What reserve would you hold if you priced with \omterm{def:dv:local-volatility:model}{local volatility}?
\item State the \emph{named result}: the gap between the local-volatility and the Heston price of the one-year monthly
\omterm{def:dv:asians-lookbacks-cliquets-and-forward-starts:cliquet}{cliquet} when both fit the same surface.
\item In one sentence: what does a \omterm{def:dv:asians-lookbacks-cliquets-and-forward-starts:cliquet}{cliquet} price that a vanilla does not?
\end{enumerate}
\end{problem}

\section{Interview questions}

\begin{interviewq}[$\star$ \iqroles{trader, researcher}]\label{iq:dv:asians-lookbacks-cliquets-and-forward-starts:1}
Why is an \omterm{def:dv:asians-lookbacks-cliquets-and-forward-starts:asian}{Asian option} cheaper than a vanilla, and by roughly how much?
\end{interviewq}

\begin{interviewq}[$\star$ \iqroles{developer}]\label{iq:dv:asians-lookbacks-cliquets-and-forward-starts:2}
How would you price an arithmetic Asian by Monte Carlo efficiently?
\end{interviewq}

\begin{interviewq}[$\star\star$ \iqroles{researcher}]\label{iq:dv:asians-lookbacks-cliquets-and-forward-starts:3}
Price a \omterm{def:dv:asians-lookbacks-cliquets-and-forward-starts:fwdstart}{forward-start option} under Black--Scholes. What does it depend on in a smile model?
\end{interviewq}

\begin{interviewq}[$\star\star$ \iqroles{trader}]\label{iq:dv:asians-lookbacks-cliquets-and-forward-starts:4}
What is a \omterm{def:dv:asians-lookbacks-cliquets-and-forward-starts:cliquet}{cliquet} exposed to, and why can two models that fit the same surface price it differently?
\end{interviewq}

\begin{interviewq}[$\star\star$ \iqroles{risk, researcher}]\label{iq:dv:asians-lookbacks-cliquets-and-forward-starts:5}
Why is \omterm{def:dv:local-volatility:model}{local volatility} a poor model for \omterm{def:dv:asians-lookbacks-cliquets-and-forward-starts:cliquet}{cliquets}?
\end{interviewq}

\begin{interviewq}[$\star\star\star$ \iqroles{trader, risk}]\label{iq:dv:asians-lookbacks-cliquets-and-forward-starts:6}
You are short a book of \omterm{def:dv:asians-lookbacks-cliquets-and-forward-starts:reverse}{reverse cliquets}. What are your risks, and how do you hedge them?
\end{interviewq}
